\begin{abstract}
Reinforcement learning (RL) can train language model agents to coordinate across multiple objectives simultaneously, but distilling these policies into smaller models via supervised fine-tuning (SFT) fails catastrophically. We present a diagnostic framework that localizes this failure across three levels of analysis, demonstrated on MEM1, a 7B RL agent for multi-objective retrieval-augmented QA. First, we establish that per-objective Exact Match degrades monotonically with objective count ($0.62 \to 0.341$ from $N\!=\!1$ to $16$), setting ``slow the slope'' as the distillation target. Second, SVD analysis reveals that RL weight changes concentrate $2$--$6\times$ supra-random along the base model's principal directions, while SFT recovers significantly less of this geometric structure (alignment ratio $1.47\times$ vs.\ $2.78\times$). Third, matched-sample DPO---a stronger signal than SFT---improves over SFT by only $+1.0$ percentage points at $N\!=\!8$ ($p = 0.163$, not significant), and we show that an earlier apparent advantage was partly a truncation confound. A causal intervention---providing oracle retrieval at inference---recovers em\_partial from $0.025$ to $0.498$ ($+47.3$pp, $p < 0.0001$), and even standard BM25 retrieval recovers $+12.2$pp, confirming that the student learns a viable search strategy but cannot execute it without the retrieval backend. A 7B capacity-isolation experiment quantifies the model-size contribution at \emph{at most} $\approx$12--14\% of the teacher--student gap, remarkably stable across all objective counts. Our findings, drawn from a single MEM1-style multi-objective RL policy, redirect attention from training-signal richness toward the structural incompatibility between RL procedural policies and single-pass distillation.
\end{abstract}
